\chapter{Segunda forma fundamental. Curvaturas}
\noindent
Sea $S$ una superficie y $p\in S$, ¿cuántos vectores $n\in \mathbb{R}^3$ cumplen $n\perp T_pS$ con $|n| = 1$? Hay exactamente dos. Si elegimos uno, orientamos $T_pS$, porque dividimos las bases $\{u,v\}$ de $T_pS$ en dos clases:
\begin{itemize}
    \item Aquellas para las que $\det(u,v,n)>0$.
    \item Aquellas para las que $\det(u,v,n)<0$.
\end{itemize}
Además, un cambio de $n$ a $-n$ invierte la orientación en $T_pS$.

\begin{definicion}
    Diremos que la superficie $S$ es \underline{orientable} si se pueden orientar todos sus planos tangentes mediante una aplicación diferenciable, es decir, si existe $N:S\to \mathbb{R}^3$ diferenciable tal que
    \begin{equation*}
        \forall p\in S \quad \text{se tiene que} \quad |N(p)| = 1 \quad \text{y}\quad   N(p)\perp T_pS
    \end{equation*}
    Si $S$ es orientable, diremos que está orientada si se ha elegido una aplicación $N$ con las características anteriores.
\end{definicion}

\begin{notacion}
    Escribiremos usualmente $N(p) = N_p$ para la aplicación $N$ anteriormente enunciada, a la que es costumbre referirnos por ``orientación'' para $S$, o también campo normal unitario o aplicación de Gau\ss.
\end{notacion}

\subsubsection{Propiedades}
\begin{enumerate}
    \item Sea $S$ una superficie conexa y orientable, entonces existen exactamente dos orientaciones sobre $S$.

        Por ser $S$ orientable, debe existir $N:S\to \mathbb{R}^3$ verificando:
        \begin{equation*}
            N_p\perp T_pS \quad \text{y}\quad |N_p| = 1 \qquad \forall q\in S
        \end{equation*}
        Fijamos dicho $N$ y supongamos que $N':S\to\mathbb{R}^3$  es otra orientación de $S$. Si hacemos el producto escalar de ambas obtenemos $\langle N,N' \rangle :S\to \mathbb{R}$, que:
        \begin{equation*}
            \langle N,N' \rangle (p) = \langle N_p,N_p' \rangle  = \pm 1
        \end{equation*}
        Como $\langle N,N' \rangle $ es continua (es diferenciable de hecho), toma valores en un discreto y $S$ es conexo tenemos que $\langle N,N' \rangle $ tiene que ser constante, con lo que hay dos posibilidades:
        \begin{itemize}
            \item Bien $\langle N,N' \rangle (p) = 1\quad \forall p\in S$, de donde $N_p = N_p'\quad \forall p\in S$.
            \item Bien $\langle N,N' \rangle (p) = -1\quad \forall p\in S$, de donde $N_p = -N_p'\quad \forall p\in S$.
        \end{itemize}
    \item Toda superficie $S$ es localmente orientable, es decir, cada punto $p\in S$ admite un entorno $V$ en $S$ de manera que $V$ es orientable\footnote{Observemos que $V$ es una superficie por ser un abierto de una superficie, por lo que tiene sentido hablar de orientablidad de $V$.}.

        Sea $X:U\to\mathbb{R}^3$ una parametrización de $S$ con $p\in X(U) = V$. Fijado $(u,v)\in U$, podemos considerar $T_{X(u,v)}S$, plano para el que $X$ proporciona una base:
        \begin{equation*}
            \{X_u(u,v), X_v(u,v)\}
        \end{equation*}
        Así, $\{X_u,X_v\}$ forman una base de $T_XS$. Definimos ahora $N^X:U\to\mathbb{R}^3$ dado por:
        \begin{equation*}
            N^X = \frac{X_u\land X_v}{|X_u\land X_v|}
        \end{equation*}
        que tiene módulo 1 y es perpendicular a $X_u$ y a $X_v$, luego a $T_XS$. Para orientar $V$ falta considerar finalmente $N:V\to\mathbb{R}^3$ como la composición $N = N^X\circ X^{-1}$, que es diferenciable. Hemos conseguido orientar $V$.

        Observemos que una vez fijemos una parametrización estamos así fijando una orientación.
\end{enumerate}

\begin{ejemplo}
    Veamos ejemplos de superficies orientables:
    \begin{enumerate}
        \item Para un plano afín $P\subset \mathbb{R}^3$  que pasa por un punto $p_0\in \mathbb{R}^3$ y tiene por vector normal $a\in \mathbb{R}^3$:
            \begin{equation*}
                P = \{p\in \mathbb{R}^3 : \langle p-p_0,a \rangle =0\}
            \end{equation*}
            Tenemos que $T_pP = \vec{P} = \langle a \rangle^\perp$. Así, vemos que $a\perp T_pP$ con $|a| =1$. Si tomamos $N:P\to \mathbb{R}^3$ dada por $N_p=a$, hemos orientado $P$.
        \item Para una esfera $S^2(p_0,R)$ con $p_0\in \mathbb{R}^3$ y $R>0$. Para $p\in S^2(p_0,R)$ teníamos:
            \begin{equation*}
                T_pS^2(p_0,R) = \{v\in \mathbb{R}^3 : v\perp p-p_0\} = \langle p-p_0 \rangle^\perp
            \end{equation*}
            definimos una aplicación de Gau\ss\ $N:S^2(p_0,R)\to \mathbb{R}^3$ dada por:
            \begin{equation*}
                N_p = \frac{p-p_0}{|p-p_0|} = \frac{1}{R}(p-p_0)
            \end{equation*}
        \item Todos los grafos son orientables, porque existe una parametrización que los cubre por completo.

            Tenemos una aplicación diferenciable $h:U\subset\mathbb{R}^2\to \mathbb{R}^3$ que define una superficie: 
            \begin{equation*}
                \Gr g = S = \{(x,y,z)\in \mathbb{R}^3 : (x,y)\in U, \ z = h(x,y)\}
            \end{equation*}
            Podemos tomar la parametrización $X:U\to \Gr h$ dada por:
            \begin{equation*}
                X(u,v) = (u,v,h(u,v))
            \end{equation*}
            que cumple $X(U)=S$. Ahora, tenemos:
            \begin{equation*}
                X_u = (1, 0, h_u), \qquad X_v = (0, 1, h_v)
            \end{equation*}
            Sea $N^X:U\to \mathbb{R}^3$ dada por
            \begin{equation*}
                N^x = \frac{(-h_u,-h_v,1)}{\sqrt{1+h_u^2+h_v^2}} = \frac{1}{\sqrt{1+|\nabla h|^2}}(-h_u,-h_v,1)
            \end{equation*}
            Con este $N^X$ definimos la correspondiente aplicación de Gau\ss, $N=N^X\circ X^{-1}$.
        \item Imágenes inversas de un valor regular. Tenemos $f:O\subset\mathbb{R}^3\to \mathbb{R}$ y $a\in \mathbb{R}$ un valor regular de $f$, con $f^{-1}(\{a\})\neq \emptyset $. Si tomamos:
            \begin{equation*}
                S = \{(x,y,z)\in \mathbb{R}^3 : f(x,y,z)=a\}
            \end{equation*}
            es una superficie. Para $(x,y,z)\in S$, tenemos que:
            \begin{equation*}
                T_{(x,y,z)}S = \left\langle \nabla f_{(x,y,z)} \right\rangle^\perp
            \end{equation*}
            Definimos así una aplicación de Gau\ss\ $N:S\to \mathbb{R}^3$ por:
            \begin{equation*}
                N(x,y,z) = \frac{\nabla f}{|\nabla f|}(x,y,z) = \frac{1}{\sqrt{f_x^2 + f_y^2 + f_z^2}} (f_x,f_y,f_z)
            \end{equation*}
        \item \textbf{Teorema de Brower-Samelson.} Toda superficie $S\subset\mathbb{R}^3$ compacta (basta con cerrada) es orientable\footnote{Notemos que en Topología II vimos ejemplos de superficies compactas no orientables, pero en dicha asignatura nuestra definición de superficie era distinta, pues admitía autointersecciones, no como en esta asignatura.}.
    \end{enumerate}
\end{ejemplo}

\begin{ejercicio}
    Un par de ejercicios:
    \begin{enumerate}
        \item Supongamos que $S = S_1\cup S_2$ para $S_i$ una superficie orientable para cada $i=1,2$. Si $S_1\cap S_2$ es conexa entonces $S$ es orientable.
        \item Supongamos que $S=S_1\cup S_2$, para $S_i$ una superficie orientable y conexa para cada $i=1,2$. Si $S_1\cap S_2$ tiene exactamente dos componentes conexas $A$ y $B$. Si $S_1$ y $S_2$ se pueden orientar de forma que las orientaciones inducidas en $A$ y en $B$ son opuestas, entonces $S$ no es orientable. Aplicarlo a la cinta de Möbius, $M = \im X$ para $X:\mathbb{R}^2\to \mathbb{R}^3$ dada por:
            \begin{equation*}
                X(u,v) = \left(\left(2-v\sen \frac{u}{2}\right)\sen u,\left(2-v\sen\frac{u}{2}\right)\cos u,v\cos\frac{u}{2}\right)
            \end{equation*}
    \end{enumerate}
\end{ejercicio}
